\begin{table}[htbp]
\centering
\caption{Dez configurações do estudo de ablação em ordenação exploratória por F1 nos respectivos testes.}
\label{tab:ablation-top10}
\small
\begin{tabular}{clccc}
\toprule
\# & Configuração & Precisão & Sensibilidade & F1 \\
\midrule
1 & Somente ML: XGBoost ($\geq$0,7) & 1,000 & 0,892 & \textbf{0,943} \\
2 & Somente ML: XGBoost ($\geq$0,5) & 0,985 & 0,892 & 0,936 \\
3 & Somente ML: XGBoost ($\geq$0,3) & 0,944 & 0,919 & 0,932 \\
4 & Somente ML: RF+SMOTE ($\geq$0,5) & 0,957 & 0,905 & 0,931 \\
5 & Consenso ML (limiar 0,7; 3 modelos) & 1,000 & 0,865 & 0,928 \\
6 & Somente ML: RF+SMOTE ($\geq$0,3) & 0,869 & 0,986 & 0,924 \\
7 & Consenso ML (limiar 0,5; 3 modelos) & 0,957 & 0,892 & 0,923 \\
8 & Híbrido OR RF+SMOTE $\geq$0,7 + Regras$\geq$9 & 0,985 & 0,865 & 0,921 \\
9 & Híbrido OR RF+SMOTE $\geq$0,7 + Regras$\geq$8 & 0,985 & 0,865 & 0,921 \\
10 & Somente ML: RF+SMOTE ($\geq$0,7) & 0,985 & 0,865 & 0,921 \\
\bottomrule
\end{tabular}
\vspace{2pt}
\par\small Nota: As linhas 1 a 3 provêm dos acréscimos com XGBoost pelo fluxo Modal; as linhas 4 a 10 provêm da ablação original, cujo consenso utiliza RF, RF+SMOTE e GB. A tabela apresenta um recorte de configurações, não as dez primeiras de todas as linhas do artefato ampliado. A ordenação por $F_1$ foi obtida após avaliação de grades de limiares nos respectivos testes internos, sem conjunto intermediário de validação para essa escolha; não constitui seleção validada nem avaliação independente dos limiares. A coincidência de proporção de teste (30\%) e semente (42) não comprova identidade das linhas. Os resultados com XGBoost aqui apresentados não são os do teste interno QW1.
\end{table}
